\documentclass{article}
\usepackage{mathtools} 
\usepackage{fontspec}
\usepackage[UTF8]{ctex}
\usepackage{amsthm}
\usepackage{mdframed}
\usepackage{xcolor}
\usepackage{amssymb}
\usepackage{amsmath}
\usepackage{tikz}
\usepackage{pgfplots}
\usepgfplotslibrary{polar}
% \pgfplotsset{compat=1.18}


% 定义新的带灰色背景的说明环境 zremark
\newmdtheoremenv[
  backgroundcolor=gray!10,
  % 边框与背景一致，边框线会消失
  linecolor=gray!10
]{zremark}{说明}

% 通用矩阵命令: \flexmatrix{矩阵名}{元素符号}{行数}{列数}
\newcommand{\flexmatrix}[4]{
  \[
  #1 = \begin{pmatrix}
    #2_{11}     & #2_{12}     & \cdots & #2_{1#4}   \\
    #2_{21}     & #2_{22}     & \cdots & #2_{2#4}   \\
    \vdots      & \vdots      & \ddots & \vdots     \\
    #2_{#31}    & #2_{#32}    & \cdots & #2_{#3#4}
  \end{pmatrix}
  \]
}

% 简化版命令（默认矩阵名为A，元素符号为a）: \quickmatrix{行数}{列数}
\newcommand{\quickmatrix}[2]{\flexmatrix{A}{a}{#1}{#2}}

\begin{document}
\title{注释 11.3}
\author{张志聪}
\maketitle

\begin{zremark}
  第九章$\S 5$习题7的瑕积分版本

  设$\int_{0}^{\pi/2} f(sin \, x) dx$为收敛的瑕积分，则有：
  \begin{itemize}
    \item (1) $\int_{0}^{\pi/2} f(sin \, x) dx = \int_{0}^{\pi/2} f(cos \, x) dx$，其中瑕点为$x = 0$；
    % \item (2) $\int_{0}^{\pi} xf(sin \, x) dx = \frac{\pi}{2} \int_{0}^{\pi} f(sin \, x) dx $，
    %       其中瑕点为$x = \pi$.
  \end{itemize}
\end{zremark}
\textbf{证明：}
\begin{itemize}
  \item (1)

        已知$\int_{0}^{\pi/2} f(sin \, x) dx$收敛，由瑕积分定义，
        \begin{align*}
          \int_{0}^{\pi/2} f(sin \, x) dx
           & = \lim \limits_{u \to 0^+} \int_{u}^{\pi/2} f(sin \, x) dx
        \end{align*}

        令$x = \frac{\pi}{2} - t, dx = -dt, x: u \to \pi/2, t: \pi/2 - u \to 0$.
        \begin{align*}
          \lim\limits_{u \to 0^+} \int_{u}^{\pi/2} f (sin x) dx
           & = \lim\limits_{u \to 0^+} \int_{\pi/2 - u}^{0} - f (sin(\frac{\pi}{2} - t)) dt \\
           & = \lim\limits_{u \to 0^+} \int_{0}^{\pi/2 - u} f (cost) dt
        \end{align*}
        令$v = \pi/2 - u, u \to 0^+, v \to (\pi/2)^-$.
        \begin{align*}
          \lim\limits_{u \to 0^+} \int_{0}^{\pi/2 - u} f (cost) dt
           & = \lim\limits_{v \to (\pi/2)^-} \int_{0}^{v} f (cost) dt \\
           & = \int_{0}^{\pi/2} f (cost) dt
        \end{align*}
        所以
        \begin{align*}
          \int_{0}^{\pi/2} f (sin x) dx
           & = \int_{0}^{\pi/2} f (cos x) dx
        \end{align*}

  % \item (2)

  %       \begin{align*}
  %         \int_{0}^{\pi} xf(sin \, x) dx
  %          & = \lim\limits_{u \to \pi^{-}} \int_{0}^{u} xf(sin \, x) dx 
  %       \end{align*}
  %       令$t = \pi - \theta, d\theta = - dt, \theta: 0 \to u, t: \pi \to \pi - u$.
  %       \begin{align*}
  %         \lim\limits_{u \to \pi^{-}} \int_{0}^{u} xf(sin \, x) dx 
  %         & = \lim\limits_{u \to \pi^{-}} \int_{\pi}^{\pi - u} - (\pi - t) f(sin \, (\pi - t)) dt \\
  %         & = \lim\limits_{u \to \pi^{-}} \int_{\pi - u}^{\pi} (\pi - t) f(sin \, t) dt 
  %       \end{align*}
  %       令$v = \pi - u, u \to \pi^-, v \to 0^+$，所以
  %       \begin{align*}
  %         \lim\limits_{u \to \pi^{-}} \int_{\pi - u}^{\pi} (\pi - t) f(sin \, t) dt 
  %         & = \lim\limits_{v \to 0^{+}} \int_{v}^{\pi} (\pi - t) f(sin \, t) dt \\
  %         & = \lim\limits_{v \to 0^{+}} \int_{v}^{\pi} \pi f(sin \, t) dt - \lim\limits_{v \to 0^{+}} \int_{v}^{\pi} t f(sin \, t) dt \\
  %       \end{align*}

\end{itemize}


\begin{zremark}
  瑕积分有
  \begin{align*}
    \int_{a}^{b} f(x) dx = A \\
    \int_{a}^{b} g(x) dx = B
  \end{align*}
  其中$x = a$是$f(x)$的瑕点，$x=b$是$g(x)$的瑕点，我们有
  \begin{align*}
    \int_{a}^{b} f(x) + g(x) dx = \int_{a}^{b} f(x) dx + \int_{a}^{b} g(x) dx = A + B
  \end{align*}
\end{zremark}

\textbf{证明：}
取$c \in (a,b)$，我们有
\begin{align*}
  \int_{a}^{b} f(x) dx = \int_{a}^{c} f(x) dx + \int_{c}^{b} f(x) dx \\
  \int_{a}^{b} g(x) dx = \int_{a}^{c} g(x) dx + \int_{c}^{b} g(x) dx
\end{align*}
于是
\begin{align*}
  \int_{a}^{b} f(x) dx + \int_{a}^{b} g(x) dx
   & = \int_{a}^{c} f(x) dx + \int_{c}^{b} f(x) dx + \int_{a}^{c} g(x) dx + \int_{c}^{b} g(x) dx
\end{align*}
考虑，因为任意$u \in [a, c]$，$g$在$[x, c]$上可积。于是定义
\begin{align*}
  G(u) = \int_{u}^{c} g(x) dx
\end{align*}
这个变下限的定积分定义的函数$G(u)$，由定理9.9可知，在$[a, c]$上函数$G(u)$连续，
所以
\begin{align*}
  \lim\limits_{u \to a^+} G(u)
   & = \lim\limits_{u \to a^+} \int_{u}^{c} g(x) dx \\
   & = G(a)                                         \\
   & = \int_{a}^{c} g(x) dx
\end{align*}
于是
\begin{align*}
  \int_{a}^{c} f(x) dx + \int_{a}^{c} g(x) dx
   & = \lim\limits_{u \to a^+} \int_{u}^{c} f(x) dx + \lim\limits_{u \to a^+} \int_{u}^{c} g(x) dx \\
   & = \lim\limits_{u \to a^+} \int_{u}^{c} f(x) dx + \int_{u}^{c} g(x) dx                         \\
   & = \lim\limits_{u \to a^+} \int_{u}^{c} f(x) + g(x) dx
\end{align*}
（注：利用了极限性质）

同理，我们有
\begin{align*}
  \int_{c}^{b} f(x) dx + \int_{c}^{b} g(x) dx
   & = \lim\limits_{u \to b^-} \int_{c}^{u} f(x) + g(x) dx
\end{align*}
综上，$x=a, x=b$是$f(x) + g(x)$的瑕点，按照瑕积分定义
\begin{align*}
  \int_{a}^{b} f(x) + g(x) dx
   & = \lim\limits_{u \to a^+} \int_{u}^{c} f(x) + g(x) dx + \lim\limits_{u \to b^-} \int_{c}^{u} f(x) + g(x) dx \\
   & = \int_{a}^{c} f(x) dx + \int_{a}^{c} g(x) dx + \int_{c}^{b} f(x) dx + \int_{c}^{b} g(x) dx                 \\
   & = \int_{a}^{b} f(x) dx + \int_{a}^{b} g(x) dx                                                               \\
   & = A + B
\end{align*}

\begin{zremark}
  偶函数的瑕积分收敛，我们有
  \begin{align*}
    \int_{-a}^a f(x) dx
     & = 2 \int_{0}^a f(x) dx
  \end{align*}
  其中，$x = a, x = -a$是瑕点。
\end{zremark}

\textbf{证明：}
\begin{align*}
  \int_{-a}^a f(x) dx
   & = \int_{-a}^0 f(x) dx + \int_{0}^a f(x) dx
\end{align*}
其中
\begin{align*}
  \int_{-a}^0 f(x) dx
   & = \lim\limits_{u \to -a^+} \int_{u}^{0} f(x) dx
\end{align*}
令$t = -x, dx = -dt, x: u \to 0, t: -u \to 0$，
\begin{align*}
  \lim\limits_{u \to -a^+} \int_{u}^{0} f(x) dx
   & = \lim\limits_{u \to -a^+} \int_{- u}^{0} - f(-t) dt \\
   & = \lim\limits_{u \to -a^+} \int_{0}^{-u} f(-t) dt
\end{align*}
令$v = -u, u: \to -a^+, v: \to a^-$，
\begin{align*}
  \lim\limits_{u \to -a^+} \int_{0}^{-u} f(-t) dt
   & = \lim\limits_{v \to a^-} \int_{0}^{v} f(-t) dt
\end{align*}
因为$f(x)$是偶函数，所以
\begin{align*}
  \lim\limits_{v \to a^-} \int_{0}^{v} f(-t) dt
   & = \lim\limits_{v \to a^-} \int_{0}^{v} f(t) dt \\
   & = \int_{0}^{a} f(t) dt
\end{align*}

综上
\begin{align*}
  \int_{-a}^a f(x) dx
   & = \int_{-a}^0 f(x) dx + \int_{0}^a f(x) dx \\
   & = \int_{0}^a f(x) dx + \int_{0}^a f(x) dx  \\
   & = 2 \int_{0}^a f(x) dx
\end{align*}

\end{document}